\section{研究现状分析}
\label{section-analysis}

根据对控制系统状态机建模与验证领域研究现状的系统梳理，结合本文的研究定位与目标，现对领域现状进行深入分析，指出关键问题，总结发展趋势，并明确本文研究的着力点。

\subsection{论文研究领域存在的问题}
\label{section-analysis-problem}

当前控制系统状态机建模与验证领域虽取得了显著进展，但在应对高复杂度系统需求与提升开发效率方面仍面临系列瓶颈问题，亟需突破性创新。

在模型构建环节，尽管基于LLM的自动化建模方法已展现出潜力，其实际效能仍受限于关键结构特征的建模缺陷。现有LLM驱动的单步建模策略普遍缺乏对状态机固有复杂结构（如层次化状态嵌套、并发区域同步、实时性约束表达）的有效引导机制，导致生成模型常出现状态迁移逻辑碎片化、并发行为描述缺失及时间属性建模不精确等问题。典型体现为：LLM虽能识别基本状态与事件，但对复合状态内的子状态行为、并行执行的子状态机间协调机制，以及状态最小/最大持续时间等关键实时约束的理解与表达存在系统性不足，需严重依赖人工审查与修正方能满足控制领域要求。这种结构表达能力的缺失，使得生成模型往往停留在概念层面，难以直接应用于高可靠性控制系统的工程实现。

在验证资产生成环节，传统方法面临覆盖度与形式化表达的双重困境。人工设计测试场景不仅效率低下，且难以全面覆盖系统关键路径、边界条件及异常情况；现有基于LLM的辅助生成技术虽能提升场景多样性，但生成的验证剖面多停留在自然语言描述层面，缺乏向形式化性质规约（如LTL/CTL公式）的系统性转化机制。尤其对控制系统至关重要的时间属性（如响应时限、事件时序约束）与安全属性（如状态互斥、故障恢复），LLM难以自动生成可被形式化验证工具直接处理的性质断言，导致验证环节仍需人工编制形式化规约，严重制约验证自动化水平。同时，验证剖面对模型关键元素的覆盖度评估机制缺失，难以确保验证的完备性。

在模型验证环节，现有技术路径存在效率与深度的矛盾。模拟测试方法虽工程友好，但受限于测试用例质量，难以发现深层次并发错误和时序违例；形式化验证方法（如模型检测）虽可提供理论完备性保证，但在处理具有复杂实时约束、多并发单元的控制系统模型时，普遍遭遇状态空间爆炸问题。现有状态缩减策略（如抽象、切片）在保留时间属性与并发交互语义方面存在局限，常需在精度与可计算性间妥协。针对LLM生成模型特有的结构特征（如深度嵌套状态），尚未形成适配的高效验证框架，制约了形式化验证在工业级复杂系统中的规模化应用。

在模型修复环节，现有方法难以有效利用验证反馈实现闭环优化。传统人工修复高度依赖专家经验，效率低且易引入新错误；基于规则的自动修复方法面对复杂逻辑变更时灵活度不足；而新兴的LLM修复技术在程序代码领域已有探索，但针对形式化状态机模型的修复仍处空白。核心问题在于：现有修复机制缺乏对验证环节输出的结构化反馈（如反例路径、性质违反点）的有效解析与转化能力，无法指导LLM精准定位缺陷根源（如特定迁移条件错误、时间约束缺失）并生成语义一致的结构化修复方案。验证与修复环节的割裂，导致难以建立高效的“设计-验证-修复”迭代闭环，阻碍模型质量的持续收敛。

\subsection{论文研究领域的发展趋势}
\label{section-analysis-trend}

面对上述挑战，本领域研究正呈现向智能化、自动化、融合化深度演进的发展趋势，旨在构建覆盖全生命周期的模型质量保障体系。

LLM与形式化方法的深度融合成为核心发展方向。LLM在自然语言理解、模式生成及上下文推理方面的超强能力，为解决传统形式化方法面临的语义鸿沟、知识获取与工程适用性问题提供新路径。研究重点正转向：1) 设计具有领域结构引导能力的多步LLM建模策略，通过分步引导与模块化生成，提升对层次化状态、并发行为等复杂结构的建模精度；2) 开发基于LLM的性质规约自动生成框架，实现从自然语言需求到形式化断言的端到端转换，尤其侧重时间逻辑与安全属性的精确表达；3) 探索LLM辅助的形式化验证加速机制，如利用LLM指导抽象精化策略或约束求解方向，缓解状态爆炸问题，提升对含复杂时间约束模型的验证效率。

验证资产的自动化与智能化生成成为效率提升的关键。重点趋势包括：1) 发展基于模型结构感知的验证剖面生成技术，结合状态迁移图分析与需求语义解析，自动生成覆盖关键路径和风险场景的高价值测试序列；2) 推动LLM驱动的多层级性质生成框架构建，利用其逻辑推理能力自动推导安全属性（如状态互斥）、活性属性（如事件响应）和实时性质（如时限满足），并输出机器可处理的形式化规约；3) 构建验证反馈驱动的迭代优化机制，利用反例数据动态完善验证剖面库与性质集，形成自增强的验证能力。

“生成-验证-修复”一体化闭环成为提升模型可靠性的主流范式。研究焦点正向全流程协同优化迁移：1) 构建验证反馈（尤其是反例）驱动的模型修复机制，结合LLM的语义理解与代码生成能力，自动定位缺陷元素（如错误迁移守卫、缺失状态）并生成语义一致的修复建议；2) 设计基于时间自动机理论的形式化修复验证框架，确保修复方案不引入新违例且满足实时性约束；3) 建立多轮迭代优化流程，通过持续验证反馈修正模型缺陷，直至达到预设的质量阈值。这种闭环范式将显著提升模型开发的收敛速度与最终可信度。

\subsection{研究现状分析结论}
\label{section-analysis-conclusion}

综合分析表明，控制系统状态机建模与验证领域正经历从孤立工具应用到智能全流程协同的范式转变。本文研究将聚焦以下关键问题的突破：

针对模型生成的结构性缺陷问题，本文需突破现有LLM单步建模局限，发展多步引导式建模方法，显式融入层次化状态、并发区域及时间约束等控制领域特征，确保生成模型具备可直接用于工程实现的完整结构与精确语义。核心在于设计融合领域知识的结构化生成策略，提升模型的结构合理性与领域一致性。

针对验证资产的形式化缺失问题，本文需构建模型元素驱动的性质自动生成框架，利用LLM的语义映射与逻辑推理能力，实现需求文本（特别是时间/安全约束）向形式化性质规约（如LTL/CTL）的系统性转换。关键在于建立性质元模型与LLM提示的协同机制，确保生成性质的可验证性与高覆盖度。

针对验证效率与深度不足问题，本文需创新基于验证剖面的混合验证方法，通过剖面引导的针对性测试缩减状态空间，结合形式化方法进行深度性质核查。重点解决时间约束的表达与验证难题，发展适配复杂嵌套结构的增量式验证策略，实现验证效率与严格性的平衡。

针对修复环节的反馈利用不足问题，本文需建立反例驱动的迭代修复机制，通过LLM精准解析验证反馈（如反例路径），定位缺陷逻辑元素（如错误迁移/守卫），生成结构保持的修复方案，并形式化验证其正确性。核心在于构建“缺陷分析-修复生成-验证确认”的闭环，实现模型质量的自动化持续提升。

综上所述，本论文研究旨在通过LLM与形式化方法的深度融合创新，系统性解决控制系统状态机在建模引导、性质生成、高效验证及反馈修复中的关键挑战，建立覆盖模型全生命周期的迭代式增强方法，为高可信复杂控制系统的开发提供新的技术路径与方法论支撑。研究结果预期在自动化建模精度、形式化验证效率及闭环迭代效能等方面实现显著突破。

